\begingroup
\makeatletter\tagsleft@true\let\veqno\@@leqno\makeatother
\clearpage\SourceLeaf{337}\ArithmeticHead{330}{L}{ANALYSIS.\quad NACHLASS.}
\begin{center}\bfseries BEMERKUNGEN.\end{center}
\par\indent Die hier zusammengestellten Aufzeichnungen stehen zerstreut zwischen astronomischen Rechnungen und Entwicklungen, mit denen aber nicht immer ein unmittelbarer Zusammenhang erkennbar ist. --- Die Funktion, von der der art. [1.] handelt, stimmt, von dem Faktor $\frac43$ abgesehen, mit der im art. 90 der \textit{Theoria motus} (1809) mit $X$ bezeichneten Funktion überein, siehe Werke VII (1906), S. 116, wo die hier angegebenen Reihen- und Kettenbruchentwicklungen auftreten, vergl. auch in der Abhandlung \textit{circa seriem} (1812), Werke III, S. 137. Das Verfahren, durch das hier der Kettenbruch aus der Reihe abgeleitet wird, ist das auf \textsc{Lambert}*) zurückgehende Divisionsverfahren, im wesentlichen also dasselbe, wie das im art. 12 der Abhandlung von 1812, Werke III, S. 134 im allgemeinen Falle angewandte. Die am Schluß des art. [1.] gegebene Entwicklung nach Potenzen von $\sin\frac12\varphi$ stimmt mit der Reihe $R$ des art. [3.] überein; wir sehen hier wieder durch das Divisionsverfahren:
\[Q-P=\alpha R,\quad R-Q=\beta S,\quad S-R=\gamma T,\quad T-S=\delta U,{.}{.}{.}\]
\noindent den Kettenbruch für $\frac\varphi{\cos\varphi}$ entstehen, wo man jetzt $s=\sin\varphi$ zu setzen hat. Dieser Kettenbruch steht auch in der Abhandlung von 1812, Werke III, S. 137, Gleich. [35]. Natürlich sind alle hier auftretenden Reihen besondere Fälle von $F(\alpha,\beta,\gamma,x)$. Das gleiche gilt von den Reihen des art. [2.]; in der Bezeichnung der Abhandlung von 1812, Werke III, S. 134, ist der erste Kettenbruch des art. [2.]:
\[\begin{aligned}
&G\left(\frac{12}2,2\ \frac{13}2,x^2\right),\\[.6em]
&\text{der zweite: }G\left(\frac{12}2,1,\frac{11}2,x^2\right),\\[.6em]
&\text{der dritte: }F\left(\frac{12}2,1,\frac{11}2,x^2\right);
\end{aligned}\]
\noindent vergl. die im Astronom. Jahrbuch für 1811, Werke VII (1906), S. 301 gegebenen Entwicklungen für $A$, ferner Werke III, S. 209. An der Hand von \textsc{Eulers} Abhandlung \textit{Specimen transformationis singularis serierum}, Nova Acta Acad. Petropol. \textbf{12} (1794) 1801, S. 58---70, fand \textsc{Gauss} auch den \textls[120]{allgemeinen Gesichtspunkt}\linebreak für die speziellen Entwicklungen der artt. [1.]---[3.]; aus dieser Abhandlung \textsc{Eulers} (\S\ 10, a. a. O. S. 61, 62) stammt das »Theorem« des art. [4.]. Die darauf folgende »Verwandlung«, die in den Bezeichnungen der Abhandlung \textit{circa seriem}
\[\begin{aligned}
G(\alpha,\gamma-\beta,\gamma,x)&=\frac{\gamma(\gamma-\alpha-\beta)}{(\gamma-\alpha)(\gamma-\beta)}\\[.6em]
&\quad+\frac{\alpha\beta}{(\gamma-\alpha)(\gamma-\beta)}G(\gamma-\alpha,\beta,\gamma,x)
\end{aligned}\]
\noindent autet, ergibt sich aus diesem \textsc{Euler}schen Theorem mit Hinzunahme der Relation (siehe Werke III, S. 133, Gl. [19])
\begin{equation}
\begin{aligned}
F(\alpha,\beta+1,\gamma+1,x)&=F(\alpha,\beta,\gamma,x)\\[.6em]
&\quad+\frac{\alpha(\gamma-\beta)}{\gamma(\gamma+1)}x\,{.}\,F(\alpha+1,\beta+1,\gamma+2,x),
\end{aligned}\tag*{($\ast$)}
\end{equation}
\noindent aus der die Kettenbruchentwicklung für den Quotienten
\[G(\alpha,\beta,\gamma,x)=\frac{F(\alpha,\beta+1,\gamma+1,x)}{F(\alpha,\beta,\gamma,x)}\]
\par\Needspace{10\baselineskip}\bigskip\noindent\rule{.25\linewidth}{.4pt}\par\smallskip
{\small\par\indent *) Siehe J. H. \textsc{Lambert}, \textit{Mémoire sur quelques propriétés remarquables des quantités transcendantes circulaires et logarithmiques}, Histoire de l’Académie, Berlin (1761) 1768, Mémoires, S. 265; siehe besonders S. 268. Während \textsc{Lambert} für die von ihm betrachteten besonderen Fälle eine Untersuchung der Konvergenz der Kettenbrüche vornimmt, fehlt eine solche Untersuchung bei \textsc{Gauss}, vergl. dazu die Bemerkung zu der weiter unten folgenden Abhandlung [III.]. Es verdient noch hervorgehoben zu werden, daß \textsc{Gauss} in späterer Zeit auf diese \textsc{Lambert}sche Abhandlung ausdrücklich Bezug genommen hat, siehe die Fußnote Werke X~2, S. 62.\par}
\clearpage\SourceLeaf{338}\ArithmeticHead{331}{R}{{\scriptsize BEMERKUNGEN ZU $F(\alpha,\beta,\gamma,x)$, ABSCHNITT I.}}
\noindent folgt. \textsc{Gauss} hat also zu der Zeit, wo er jene »Verwandlung« ausführte, die Relation ($\ast$) und damit den Schlüssel für die ganze Theorie der Reihe $F(\alpha,\beta,\gamma,x)$ gekannt. Diese Zeit läßt sich mit ziemlicher Genauigkeit bestimmen. Auf S. 43---45 des Handbuchs 18, Bd, beginnt nämlich die Untersuchung über die Identität des Kometen von 1772 mit dem von 1805, indem dort die Beobachtungen des Kometen von 1772 neu reduziert und berechnet werden. Über diese Rechnung berichtet \textsc{Gauss} an \textsc{Olbers} am 3. Januar 1806, siehe \textit{Wilhelm Olbers, Sein Leben und seine Werke} II, 1, 1900, S. 279. Andererseits finden sich auf S. 64ff. des Handbuchs Reduktionen von Beobachtungen der Pallas und Juno aus dem Februar 1806. Damit sind die Aufzeichnungen der artt. [1.]---[3.] und das »Theorem« von [4.] auf den Anfang des Jahres 1806 festgelegt. Aber auch die »Verwandlung« stammt aus dem ersten Drittel dieses Jahres, da die auf S. 147 des Handbuchs stehende algebraische Aufzeichnung (abgedruckt oben S. 112, art. [3.]) nach der Aufzeichnung Nr. 128 des \textit{Tagebuchs} aus den Monaten April---Mai 1806 herrührt.
\begin{flushright}\textsc{Schlesinger.}\end{flushright}
\par\bigskip\begin{center}\rule{.25\linewidth}{.4pt}\end{center}
\clearpage
\endgroup
